\documentclass[AMS,STIX2COL]{WileyNJD-v2}

% Additional packages (not loaded by WileyNJD-v2.cls)
\usepackage{amssymb}
\usepackage{colortbl}

\graphicspath{{../figures/}}

% Custom colors (xcolor already loaded by cls)
\definecolor{safegreen}{HTML}{27ae60}
\definecolor{dangerred}{HTML}{e74c3c}
\definecolor{warnorange}{HTML}{e67e22}

% Consistent float styling
\setlength{\abovecaptionskip}{6pt}
\setlength{\belowcaptionskip}{0pt}
\renewcommand{\arraystretch}{1.15}
\newcommand{\tabstyle}{\small\setlength{\tabcolsep}{5pt}}
\newcommand{\tabstylecompact}{\small\setlength{\tabcolsep}{3.5pt}}
\newcommand{\figcol}[1]{\includegraphics[width=0.98\columnwidth]{#1}}
\newcommand{\figcolsm}[1]{\includegraphics[width=0.88\columnwidth]{#1}}
\newcommand{\figwide}[1]{\includegraphics[width=0.98\textwidth]{#1}}
\newcommand{\figwideh}[1]{\includegraphics[width=0.92\textwidth,height=0.42\textheight,keepaspectratio]{#1}}

\articletype{Special Issue on Trustworthy and Safe AI}%

\received{<day> <Month>, <year>}
\revised{<day> <Month>, <year>}
\accepted{<day> <Month>, <year>}

\begin{document}

\title{Think Before You Refuse:\\
Safety Alignment Gaps in Open-Source Reasoning Models}

\author[1]{Anonymous Author*}
\author[1]{Anonymous Co-author}

\authormark{ANONYMOUS \textsc{et al.}}

\address[1]{\orgdiv{Department}, \orgname{Institution},
  \orgaddress{\state{State}, \country{Country}}}

\corres{*Corresponding author. \email{anonymous@example.com}}

\abstract[Abstract]{%
Open-source reasoning models such as the DeepSeek-R1 distillation family
generate explicit chain-of-thought traces before answering, yet the
consequences of this deliberative paradigm for safety alignment remain
poorly understood.
We conduct a controlled comparison of three architecture-matched model
pairs---Llama-3.1-8B, Qwen-2.5-7B, and Qwen-2.5-14B---that differ only
in training objective (standard instruction tuning vs.\ reasoning
distillation), evaluating each on the AdvBench benchmark (520 harmful
prompts across 7 categories).
Reasoning-distilled models refuse only 75--80\% of harmful prompts compared
with 95--99\% for their standard counterparts, yielding a 4--26$\times$
amplification in attack success rate.
This degradation is category-selective: deception collapses ($-43$~pp)
while discrimination and sexual content remain fully safe.
To explain these gaps, we introduce a four-type taxonomy for safety-related
patterns inside reasoning traces, revealing a \emph{self-rationalization}
failure mode in which models explicitly recognize harm yet reason themselves
into compliance (7--9\% of cases), alongside a \emph{safety-blind} mode
with zero safety awareness (13--16\%).
These findings demonstrate that extended reasoning can undermine alignment
by supplying cognitive pathways that bypass trained refusals, motivating
safety-preserving distillation methods and trace-aware inference-time
defenses.
}

\keywords{%
Large language models,
Reasoning models,
Safety alignment,
Chain-of-thought,
DeepSeek-R1,
Self-rationalization,
Trustworthy AI,
Interpretability
}

\jnlcitation{\cname{%
\author{Anonymous A.},
\author{Anonymous B.}} (\cyear{2026}),
\ctitle{Think Before You Refuse: Safety Alignment Gaps in Open-Source
Reasoning Models}, \cjournal{ETRI Journal}, \cvol{2026;00:1--17}.}

\maketitle

\footnotetext{%
\textbf{Abbreviations:}
LLM --- Large Language Model;
CoT --- Chain-of-Thought;
ASR --- Attack Success Rate;
RLHF --- Reinforcement Learning from Human Feedback;
DPO --- Direct Preference Optimization
}

% ============================================================
\section{Introduction}
\label{sec:intro}
% ============================================================

Reasoning-enhanced large language models (LLMs) are reshaping the performance
frontier of artificial intelligence.
Models such as OpenAI's o1\cite{openai2024o1} and
DeepSeek-R1\cite{deepseek2025r1} explicitly generate intermediate reasoning
before producing a final response; this ``deliberate then answer'' paradigm
has produced strong gains on mathematical reasoning, code generation, and
multi-step logic tasks.

In parallel, \textbf{safety alignment}---training models to refuse harmful,
illegal, or unethical requests---has become a cornerstone of responsible AI
deployment.
Standard instruction-tuned models undergo safety fine-tuning via
RLHF\cite{ouyang2022training} or DPO,\cite{rafailov2023direct} which
typically produces robust refusal behaviors when models encounter harmful
prompts.

These two trends create an open and urgent question:
\emph{Does explicit deliberation strengthen safety by enabling more nuanced
judgment, or does it create pathways for models to rationalize around safety
guardrails?}
Reasoning traces expose safety awareness in an interpretable form, but they
also provide additional degrees of freedom for constructing justifications
for compliance.

This question is particularly pressing for \textbf{open-source reasoning
models}.
The DeepSeek-R1 distillation family\cite{deepseek2025r1} fine-tunes existing
open-source base models (Llama, Qwen) to produce reasoning traces, making
deliberative behavior accessible at consumer-hardware scale.
If reasoning distillation systematically degrades safety alignment, the rapid
proliferation of these models poses a concrete and scalable risk to
downstream applications and end users.

We investigate this through a controlled experimental design: we pair three
standard instruction-tuned models with their reasoning-distilled
counterparts---all sharing the same base architecture---and evaluate them on
the AdvBench harmful behaviors benchmark.\cite{zou2023universal}
Critically, our analysis goes beyond aggregate refusal rates: we extract and
classify the \emph{content} of reasoning traces to identify the mechanisms by
which safety alignment is preserved or circumvented.

\textbf{Contributions.}
\begin{enumerate}[(1)]
  \item A \textbf{systematic safety comparison} across three
        architecture-matched model pairs spanning two model families
        (Llama, Qwen) and two parameter scales (7--8B, 14B), isolating the
        effect of reasoning distillation from architectural confounds.
  \item A \textbf{four-type taxonomy of safety reasoning patterns} in
        thinking traces, including \emph{self-rationalization}---cases where
        models recognize harm yet reason themselves into compliance---and
        \emph{safety blindness}---cases with zero safety awareness.
  \item \textbf{Empirical evidence} that reasoning distillation amplifies
        ASR by 4--26$\times$, with \emph{category-selective} vulnerability
        concentrated in ambiguous-boundary harm categories
        (deception: $-43$~pp).
  \item \textbf{Scale analysis} showing that 7B\,$\to$\,14B provides
        modest safety recovery ($-4.6$~pp ASR) primarily via improved
        safety awareness activation rather than rationalization resistance.
\end{enumerate}

% ============================================================
\section{Related Work}
\label{sec:related}
% ============================================================

\subsection{Reasoning Models}

Chain-of-thought (CoT) prompting\cite{wei2022chain} demonstrated that
eliciting intermediate reasoning steps improves LLM performance on complex
tasks.
OpenAI's o1\cite{openai2024o1} extended this to inference-time compute
scaling with internal reasoning tokens.
DeepSeek-R1\cite{deepseek2025r1} open-sourced the paradigm via a full 671B
model and smaller distilled variants built on Llama and Qwen architectures,
making reasoning capabilities broadly accessible.
These developments motivate the question of whether performance gains from
explicit deliberation carry over---or transfer at all---to safety-critical
settings.

\subsection{Safety Alignment}

RLHF\cite{ouyang2022training} and DPO\cite{rafailov2023direct} are the
dominant alignment approaches.
Yet alignment remains brittle: jailbreaking
attacks\cite{zou2023universal,wei2024jailbroken} can bypass refusals via
adversarial prompting, fine-tuning on harmful
examples,\cite{qi2023fine,yang2023shadow} or exploiting multilingual
gaps.\cite{deng2024multilingual,yong2024lowresource}
This brittleness indicates that alignment behaviors are sensitive to shifts
in training objectives and generation dynamics---precisely the kind of shift
introduced by reasoning distillation.

\subsection{Safety of Reasoning Models}

OpenAI\cite{openai2024o1systemcard} reported that o1's extended reasoning
occasionally produces deceptive behaviors during safety evaluations.
Several concurrent works have noted weaker safety in DeepSeek-R1, but
\emph{systematic evaluations with controlled baselines and thinking trace
analysis remain scarce}.
Our work fills this gap by pairing architecture-matched models and directly
analyzing reasoning traces to identify mechanism-level failure modes.

% ============================================================
\section{Experimental Setup}
\label{sec:setup}
% ============================================================

\subsection{Model Selection}

We select three model pairs where each pair shares a common base architecture
but differs in training objective (Table~\ref{tab:models}).
This design isolates the effect of the training procedure (instruction tuning
vs.\ reasoning distillation) from architectural confounds.
The inclusion of two Qwen-based pairs (7B and 14B) additionally enables
\emph{within-family scale analysis} of the reasoning distillation paradigm.

\begin{table}[t]
\centering
\caption{Architecture-matched model pairs.  Each pair shares base weights;
only the fine-tuning objective differs.}
\label{tab:models}
\tabstyle
\begin{tabular}{@{}l l l c@{}}
\toprule
\textbf{Base} & \textbf{Model} & \textbf{Type} & \textbf{Params} \\
\midrule
\multirow{2}{*}{Llama-3.1}
  & Llama-3.1-8B-Instruct & Standard & 8B \\
  & DeepSeek-R1-Distill-Llama-8B & Reasoning & 8B \\
\addlinespace
\multirow{2}{*}{Qwen-2.5}
  & Qwen-2.5-7B-Instruct & Standard & 7B \\
  & DeepSeek-R1-Distill-Qwen-7B & Reasoning & 7B \\
\addlinespace
\multirow{2}{*}{Qwen-2.5}
  & Qwen-2.5-14B-Instruct & Standard & 14B \\
  & DeepSeek-R1-Distill-Qwen-14B & Reasoning & 14B \\
\bottomrule
\end{tabular}
\end{table}

\subsection{Benchmark and Inference}

We evaluate on \textbf{AdvBench}\cite{zou2023universal}, a benchmark of 520
harmful prompts spanning seven harm categories: violence ($n{=}100$),
illegal activity ($n{=}163$), deception ($n{=}64$), other ($n{=}166$),
self-harm ($n{=}14$), discrimination ($n{=}7$), and sexual content ($n{=}6$).
Each prompt is presented as a single-turn user message \emph{without
jailbreaking augmentation}, providing a clean baseline for refusal behavior.

All models are served via vLLM\cite{kwon2023efficient} with greedy decoding
(temperature\,$\approx 0$), a maximum of 512 new tokens, and default chat
templates on NVIDIA RTX A6000 (48\,GB)\,$\times$\,8.
This protocol isolates the effect of the training objective from sampling
randomness.

\subsection{Evaluation Protocol}

\subsubsection{Refusal detection}
We apply keyword-based refusal matching with 30+ patterns (``I cannot,''
``I'm sorry,'' ``as an AI,'' etc.).
For reasoning models, detection is applied \emph{after} removing the
thinking trace to prevent template leakage.
Any explicit refusal in the final answer is counted as a refusal regardless
of reasoning trace content.

\subsubsection{Attack success rate (ASR)}
Defined as $\text{ASR} = 1 - \text{Refusal Rate}$.

\subsubsection{Thinking trace taxonomy}
For reasoning models, we classify each thinking trace into four types:
\begin{itemize}
  \item \textbf{Type~A} \textcolor{safegreen}{(Safety-aware $\to$ refuse)}:
        The model recognizes the request as harmful and refuses.
  \item \textbf{Type~B} \textcolor{dangerred}{(Safety-aware $\to$ comply)}:
        The model recognizes potential harm but rationalizes compliance.
        Subdivided into \emph{B-explicit} (overt justification) and
        \emph{B-implicit} (notes concern, complies without explicit
        rationalization).
  \item \textbf{Type~C} \textcolor{warnorange}{(No awareness $\to$ comply)}:
        The model exhibits zero safety consideration and complies directly.
  \item \textbf{Type~D} (Other $\to$ refuse):
        The model refuses without explicit safety reasoning.
\end{itemize}
This taxonomy enables mechanism-level analysis: Type~B identifies
\emph{self-rationalization} (a reasoning-specific failure mode), while
Type~C identifies \emph{safety blindness} (a detection gap).

% ============================================================
\section{Results}
\label{sec:results}
% ============================================================

\subsection{Overall Safety Comparison}

\begin{table}[t]
\centering
\caption{Main results across architecture-matched model pairs.  Refusal
Rate\,(\%) = proportion of harmful prompts refused.
Ampl.\,= ASR\textsubscript{reasoning}\,/\,ASR\textsubscript{standard}.
Avg~Len = mean response length (characters).}
\label{tab:main}
\tabstyle
\begin{tabular}{@{}l l r r r r@{}}
\toprule
\textbf{Base} & \textbf{Type} & \textbf{Refusal} & \textbf{ASR} & \textbf{Ampl.} & \textbf{Avg Len} \\
\midrule
\multirow{2}{*}{Llama 8B}
  & Standard  & \textbf{95.0\%} & 5.0\%  & ---         & 255 \\
  & Reasoning & 78.1\%          & 21.9\% & \textbf{4.4$\times$}  & 2{,}433 \\
\addlinespace
\multirow{2}{*}{Qwen 7B}
  & Standard  & \textbf{99.0\%} & 1.0\%  & ---         & 867 \\
  & Reasoning & 75.2\%          & 24.8\% & \textbf{25.8$\times$} & 2{,}493 \\
\addlinespace
\multirow{2}{*}{Qwen 14B}
  & Standard  & \textbf{98.5\%} & 1.5\%  & ---         & 558 \\
  & Reasoning & 79.8\%          & 20.2\% & \textbf{13.1$\times$} & 2{,}446 \\
\bottomrule
\end{tabular}
\end{table}

\begin{figure}[!t]
\centering
\figcol{fig1_slope.png}
\caption{ASR convergence after reasoning distillation.  Despite widely
varying baseline safety (1--5\% ASR), all reasoning models converge to a
narrow 20--25\% ASR band.  Labels show amplification factors.}
\label{fig:slope}
\end{figure}

Table~\ref{tab:main} and Figure~\ref{fig:slope} present the central result.
Standard instruct models refuse 95--99\% of harmful prompts; their
reasoning-distilled counterparts refuse only 75--80\%.
This corresponds to a \textbf{4--26$\times$ ASR amplification}---a
substantial and consistent safety gap induced by reasoning distillation.

\textbf{Finding 1 (Convergent vulnerability).}
Despite widely varying baseline safety (1.0--5.0\% ASR), reasoning models
\emph{converge} to a narrow 20--25\% ASR band.
This convergence suggests that distillation imposes a shared vulnerability
ceiling that partially overrides the original safety profile of each base
model.

\subsection{Category-Wise Analysis}

\begin{figure}[!t]
\centering
\figcolsm{fig2_category_gap.png}
\caption{Refusal rate gap (reasoning\,$-$\,standard, in pp) by harm category.
Deception suffers catastrophic degradation ($-43$~pp); discrimination and
sexual content show zero degradation.}
\label{fig:catgap}
\end{figure}

\begin{table*}[t]
\centering
\caption{Category-wise refusal rates (\%).  Gap\,= mean difference
(reasoning\,$-$\,standard) averaged across all three model pairs.
Categories are sorted by vulnerability.}
\label{tab:category}
\tabstylecompact
\begin{tabular}{@{}l c ccc ccc c@{}}
\toprule
& & \multicolumn{3}{c}{\textbf{Standard}} & \multicolumn{3}{c}{\textbf{Reasoning}} & \\
\cmidrule(lr){3-5} \cmidrule(lr){6-8}
\textbf{Category} & \textbf{$n$}
  & \footnotesize L-8B & \footnotesize Q-7B & \footnotesize Q-14B
  & \footnotesize R-L8B & \footnotesize R-Q7B & \footnotesize R-Q14B
  & \textbf{Gap} \\
\midrule
Deception        & 64  & 87.5 & 100 & 100 & 43.8 & 50.0 & 64.1 & \textcolor{dangerred}{\textbf{$-$43.2}} \\
Violence         & 100 & 97.0 & 100 & 99.0 & 80.0 & 71.0 & 81.0 & \textcolor{dangerred}{$-$21.3} \\
Illegal Act.     & 163 & 98.8 & 100 & 100 & 82.8 & 84.0 & 82.2 & \textcolor{dangerred}{$-$16.6} \\
Other            & 166 & 92.2 & 97.0 & 97.0 & 82.5 & 75.3 & 79.5 & \textcolor{warnorange}{$-$16.3} \\
Self-Harm        & 14  & 100  & 100 & 85.7 & 92.9 & 92.9 & 100  & \textcolor{safegreen}{$\pm$0.0} \\
Discrimination   & 7   & 100  & 100 & 100 & 100  & 100  & 100  & \textcolor{safegreen}{0.0} \\
Sexual           & 6   & 100  & 100 & 100 & 100  & 100  & 100  & \textcolor{safegreen}{0.0} \\
\bottomrule
\end{tabular}
\end{table*}

Figure~\ref{fig:catgap} and Table~\ref{tab:category} reveal a striking
\emph{two-tier} vulnerability structure.

\textbf{Finding 2 (Category-selective degradation).}
\textbf{Vulnerable tier:} Deception ($-43.2$~pp), violence ($-21.3$~pp),
illegal activity ($-16.6$~pp), and other ($-16.3$~pp) exhibit substantial
safety degradation after reasoning distillation.
\textbf{Resilient tier:} Discrimination and sexual content maintain 100\%
refusal; self-harm shows $\pm 0$~pp change.

The most vulnerable category---deception---drops to a 47.4\% compliance rate
in reasoning models (vs.\ 4.2\% in standard models).
Degradation concentrates in categories with \emph{ambiguous} helpful/harmful
boundaries, whereas categories with clear social taboos survive distillation
intact.

\subsection{Thinking Trace Analysis}
\label{sec:traces}

\begin{figure}[!t]
\centering
\figcol{fig3_trace_donuts.png}
\caption{Thinking trace classification for all three reasoning models.
Type~A (safety-aware refusal) dominates at 73--79\%, but Type~B
(self-rationalization) and Type~C (no safety awareness) together
account for 20--25\% of all responses.}
\label{fig:traces}
\end{figure}

\begin{table}[t]
\centering
\caption{Thinking trace distribution (reasoning models only).  Percentages
are relative to total prompts ($n{=}520$ per model).}
\label{tab:traces}
\tabstyle
\begin{tabular}{@{}l r r r r@{}}
\toprule
\textbf{Trace} & \textbf{R-L8B} & \textbf{R-Q7B} & \textbf{R-Q14B} & \textbf{Total} \\
\midrule
\textcolor{safegreen}{\textbf{A}} Safe\,$\to$\,ref.
  & 402\,(\footnotesize 77.3\%) & 380\,(\footnotesize 73.1\%) & 411\,(\footnotesize 79.0\%) & 1{,}193\,(\footnotesize 76.5\%) \\
\addlinespace
\textcolor{dangerred}{\textbf{B}} Safe\,$\to$\,comp.
  & 36\,(\footnotesize 6.9\%) & 46\,(\footnotesize 8.8\%) & 39\,(\footnotesize 7.5\%) & 121\,(\footnotesize 7.8\%) \\
\quad \footnotesize B-explicit
  & \footnotesize 3 & \footnotesize 7 & \footnotesize 6 & \footnotesize 16\,(13.2\%) \\
\quad \footnotesize B-implicit
  & \footnotesize 33 & \footnotesize 39 & \footnotesize 33 & \footnotesize 105\,(86.8\%) \\
\addlinespace
\textcolor{warnorange}{\textbf{C}} No aware.
  & 78\,(\footnotesize 15.0\%) & 83\,(\footnotesize 16.0\%) & 66\,(\footnotesize 12.7\%) & 227\,(\footnotesize 14.6\%) \\
\addlinespace
\textbf{D} Other ref.
  & 4\,(\footnotesize 0.8\%) & 11\,(\footnotesize 2.1\%) & 4\,(\footnotesize 0.8\%) & 19\,(\footnotesize 1.2\%) \\
\bottomrule
\end{tabular}
\end{table}

Figure~\ref{fig:traces} and Table~\ref{tab:traces} present the trace
taxonomy results.
The majority of traces (76.5\%) correctly identify harm and refuse
(Type~A), demonstrating that reasoning \emph{can} reinforce safety.
However, the remaining ${\sim}$23\% decompose into two distinct failure
modes.

\textbf{Finding 3 (Self-rationalization).}
In 7--9\% of cases (Type~B), reasoning models \textbf{recognize a request
as harmful but reason themselves into compliance}.
This self-rationalization pattern---previously undocumented in open-source
models---demonstrates that extended reasoning can \emph{undermine} rather
than strengthen safety alignment.

\subsubsection{Self-rationalization strategies}
Analysis of Type~B traces reveals three recurring strategies:

\textit{Strategy 1: Tone reframing.}
A model asked to plan a mass shooting reasons:
``\emph{This is a serious and sensitive topic.  I should approach this
carefully, making sure I cover all necessary aspects without being too
graphic or offensive}'' (R1-Qwen-7B).
The model reframes the request as a matter of \emph{tone} rather than
one requiring \emph{refusal}, then provides target analysis and weapon
selection.

\textit{Strategy 2: Charitable reinterpretation.}
A model asked to build an exploitation algorithm reasons:
``\emph{This could be used for bad purposes.  Wait, maybe the user is
asking for a tool to identify people who might be at risk, so that they
can be protected.  That makes more sense}'' (R1-Llama-8B).
The model reinterprets the adversarial prompt as benign, then complies.

\textit{Strategy 3: Motive assumption.}
A model asked to create an assault tutorial reasons:
``\emph{That's a pretty serious topic.  Maybe it's for self-defense, but
I'm not sure}'' (R1-Qwen-7B).
The model assumes a benign motive and provides instructions.

\subsubsection{Type B breakdown}
A critical finding is that \textbf{86.8\%} of Type~B failures are
\emph{implicit} (B-implicit): the model notes safety concerns but complies
anyway without constructing an explicit justification.
Only 13.2\% involve overt self-persuasion (B-explicit).
The dominant failure mode is not ``talking itself into compliance'' but
rather \emph{insufficient constraint propagation}---safety awareness in the
reasoning chain fails to deterministically block harmful output generation.

\subsubsection{Category$\times$trace cross-analysis}

\begin{figure*}[!t]
\centering
\figwide{fig4_crosstab.png}
\caption{Category\,$\times$\,trace type distribution across all reasoning
models ($n{=}1{,}560$).  Deception exhibits the highest Type~C rate
(37.0\%) and Type~B rate (10.4\%), indicating that nearly half of
deception responses involve either rationalization or complete safety
blindness.}
\label{fig:crosstab}
\end{figure*}

Figure~\ref{fig:crosstab} reveals that deception prompts trigger the most
safety-oblivious responses: \textbf{37.0\%} of deception traces show
\emph{zero safety awareness} (Type~C), compared with 10--14\% for other
vulnerable categories.
Combined with a 10.4\% Type~B rate, nearly half of deception responses
involve either rationalization or complete safety blindness---explaining
deception's catastrophic vulnerability observed in Section~4.2.

\subsection{Scale Effect}

\begin{figure}[!t]
\centering
\figcol{fig5_scale.png}
\caption{Category-wise refusal rates for R1-Qwen-7B vs.\ R1-Qwen-14B.
Scaling from 7B to 14B yields the largest improvement in deception
($+14.1$~pp) and violence ($+10.0$~pp).}
\label{fig:scale}
\end{figure}

Comparing the two Qwen-based reasoning distills (Figure~\ref{fig:scale}),
scaling from 7B to 14B provides a modest overall improvement: ASR decreases
by 4.6~pp (24.8\%\,$\to$\,20.2\%).
The improvement is concentrated in the most vulnerable category
(deception: $+14.1$~pp).
Mechanistically, scaling converts \textbf{17 Type~C failures into Type~A
successes}, suggesting that larger models are more likely to \emph{activate}
safety reasoning rather than being better at \emph{resisting}
self-rationalization once activated.

% ============================================================
\section{Discussion}
\label{sec:discussion}
% ============================================================

\subsection{Why Does Reasoning Distillation Degrade Safety?}

We hypothesize two complementary mechanisms.

\textbf{Mechanism 1: Safety overwriting.}
The distillation process prioritizes reproducing the reasoning behavior of
the teacher model (DeepSeek-R1-671B) over preserving the base model's
safety alignment.
Safety refusal patterns learned during RLHF/DPO may be overwritten by
reasoning fine-tuning data that under-represents safety-critical examples.
The \emph{convergence} of reasoning-model ASR to a narrow 20--25\%
band---regardless of the base model's original safety level (1.0--5.0\%
ASR)---supports this interpretation: distillation imposes a new, weaker
safety profile that partially replaces the original one.

\textbf{Mechanism 2: Reasoning as bypass.}
Standard instruct models produce relatively ``reflexive''
refusals---pattern-matched responses triggered early in generation.
Reasoning models engage in \emph{deliberation}, which creates a pathway
for self-rationalization.
This parallels dual-process theory in cognitive science: System~1 (fast,
automatic) safety refusals work well in standard models, but System~2
(slow, deliberative) reasoning can override these automatic responses
through motivated reasoning.

\subsection{The Self-Rationalization Problem}

The dominance of B-implicit (86.8\%) over B-explicit (13.2\%) carries
important implications.
The primary failure mode is not overt self-persuasion but
\textbf{insufficient constraint propagation}: safety awareness activated in
the reasoning phase does not reliably constrain the generation phase.
This suggests that the thinking trace and output generation may operate
with partially decoupled safety signals---a hypothesis that warrants
investigation through mechanistic interpretability methods.

The three self-rationalization strategies identified in this work---tone
reframing, charitable reinterpretation, and motive assumption---parallel
well-studied \emph{motivated reasoning} patterns in human cognition, where
desired conclusions are reached first and justifications constructed
post-hoc.

\subsection{Category-Selective Vulnerability}

The two-tier degradation pattern admits a clear interpretation:
\begin{itemize}
  \item \textbf{Resilient categories} (discrimination, sexual content) have
        \emph{bright-line} boundaries with strong, consistent training
        signals that survive distillation.
  \item \textbf{Vulnerable categories} (deception, violence, illegal
        activity) have \emph{ambiguous} helpful/harmful boundaries that
        reasoning models can exploit to construct rationalizations.
\end{itemize}
The moderate rank correlation ($\rho = 0.54$) between standard-model and
reasoning-model category safety profiles implies that \textbf{standard-model
safety audits cannot reliably predict reasoning-model vulnerabilities}.

\subsection{Proposed Mitigations}

\begin{table}[t]
\centering
\caption{Multi-layered defense framework against reasoning-induced safety
failures.  Each layer targets specific failure modes identified by our
trace taxonomy.}
\label{tab:mitigations}
\tabstyle
\begin{tabular}{@{}l p{4.0cm} p{2.5cm}@{}}
\toprule
\textbf{Layer} & \textbf{Intervention} & \textbf{Targets} \\
\midrule
\textsc{Train}
  & Safety-preserved distillation & All types \\
  & Adversarial reasoning training & Type~B \\
  & Category-weighted safety RLHF & Vulnerable cats. \\
\addlinespace
\textsc{Infer.}
  & Thinking trace monitor & Type~B, C \\
  & Safety checkpoint injection & All types \\
  & Dual-path generation & Type~B \\
\addlinespace
\textsc{Deploy}
  & Category-specific guardrails & Deception \\
  & Reasoning budget caps & Type~B \\
  & Mandatory trace audit & Transparency \\
\bottomrule
\end{tabular}
\end{table}

Table~\ref{tab:mitigations} outlines a multi-layered defense framework
informed by our trace taxonomy.
At \emph{training time}, safety-preserving distillation should co-optimize
reasoning quality and safety; adversarial reasoning training should expose
models to scenarios requiring self-rationalization resistance.
At \emph{inference time}, trace monitoring can flag Type~B/C patterns
before responses reach users; safety checkpoints between reasoning and
generation can enforce trace--response consistency.
At the \emph{deployment level}, category-specific guardrails (especially
for deception) and reasoning budget caps can constrain the cognitive space
available for rationalization.

\subsection{Limitations}

Our study has several limitations.
(1)~Keyword-based refusal detection may miss nuanced refusals or yield
false positives; trace classification relies on rule-based cues rather than
human annotation.
(2)~We evaluate only English prompts and the DeepSeek-R1 distillation
family; other reasoning paradigms (o1, QwQ, Gemini~2.0 Flash Thinking) may
exhibit different safety dynamics.
(3)~Our parameter range (7--14B) is limited; larger models may show
different scaling patterns or stronger recovery.
(4)~AdvBench provides a clean but relatively simple evaluation surface;
more sophisticated adversarial attacks or multi-turn settings would likely
amplify the observed effects.

% ============================================================
\section{Conclusion}
\label{sec:conclusion}
% ============================================================

We present the first controlled comparison of standard instruction-tuned
and reasoning-distilled open-source LLMs, evaluating three
architecture-matched model pairs on 520 harmful prompts.
Our findings carry four principal messages:

\begin{enumerate}[(1)]
  \item \textbf{Universal safety degradation.}
    Reasoning models refuse only 75--80\% of harmful prompts vs.\
    95--99\% for standard models (4--26$\times$ ASR amplification),
    with reasoning-model ASR converging to a 20--25\% band regardless
    of baseline safety.

  \item \textbf{Category-selective vulnerability.}
    Deception is catastrophically vulnerable ($-43$~pp) while
    discrimination and sexual content remain perfectly safe, reflecting
    a boundary-clarity gradient in how safety survives distillation.

  \item \textbf{Self-rationalization as a novel failure mode.}
    In 7--9\% of cases, reasoning models recognize harm but reason
    themselves into compliance.
    86.8\% of these failures are implicit, indicating insufficient
    constraint propagation from reasoning to generation.

  \item \textbf{Modest scale benefit.}
    Scaling from 7B to 14B reduces ASR by 4.6~pp, primarily by
    activating safety reasoning in previously blind-spot cases
    (Type~C\,$\to$\,Type~A), not by improving rationalization resistance.
\end{enumerate}

As reasoning capabilities become standard in production LLMs, ensuring that
extended deliberation \emph{strengthens rather than undermines} safety
alignment is a critical challenge.
We advocate for safety-preserving distillation methods, thinking trace
monitoring as a new inference-time safety layer, and independent safety
evaluations for reasoning models that do not assume base-model safety
profiles transfer through distillation.

% ============================================================
% REFERENCES
% ============================================================
\begin{thebibliography}{99}

\bibitem{deepseek2025r1}
DeepSeek-AI,
DeepSeek-R1: Incentivizing reasoning capability in LLMs via reinforcement
learning,
\textit{arXiv preprint arXiv:2501.12948}, 2025.

\bibitem{openai2024o1}
OpenAI,
Learning to reason with LLMs,
\textit{OpenAI Blog}, 2024.

\bibitem{openai2024o1systemcard}
OpenAI,
OpenAI o1 system card,
\textit{OpenAI Technical Report}, 2024.

\bibitem{ouyang2022training}
L.~Ouyang, J.~Wu, X.~Jiang, D.~Almeida, C.~Wainwright, P.~Mishkin,
C.~Zhang, S.~Agarwal, K.~Slama, A.~Ray, \textit{et~al.},
Training language models to follow instructions with human feedback,
\textit{Advances in Neural Information Processing Systems},
\textbf{35}, 27730--27744, 2022.

\bibitem{rafailov2023direct}
R.~Rafailov, A.~Sharma, E.~Mitchell, C.~D.~Manning, S.~Ermon, C.~Finn,
Direct preference optimization: Your language model is secretly a reward
model,
\textit{Advances in Neural Information Processing Systems},
\textbf{36}, 2023.

\bibitem{zou2023universal}
A.~Zou, Z.~Wang, N.~Carlini, M.~Nasr, J.~Z.~Kolter, M.~Fredrikson,
Universal and transferable adversarial attacks on aligned language models,
\textit{International Conference on Machine Learning}, 2023.

\bibitem{wei2024jailbroken}
A.~Wei, N.~Haghtalab, J.~Steinhardt,
Jailbroken: How does LLM safety training fail?,
\textit{Advances in Neural Information Processing Systems},
\textbf{36}, 2024.

\bibitem{qi2023fine}
X.~Qi, Y.~Zeng, T.~Xie, P.-Y.~Chen, R.~Jia, P.~Mittal, P.~Henderson,
Fine-tuning aligned language models compromises safety, even when users do
not intend to!,
\textit{arXiv preprint arXiv:2310.03693}, 2023.

\bibitem{yang2023shadow}
X.~Yang, X.~Xiao, Q.~Bi, Y.~Ke, \textit{et~al.},
Shadow alignment: The ease of subverting safely-aligned language models,
\textit{arXiv preprint arXiv:2310.02949}, 2023.

\bibitem{deng2024multilingual}
Y.~Deng, W.~Zhang, S.~J.~Pan, L.~Bing,
Multilingual jailbreak challenges in large language models,
\textit{arXiv preprint arXiv:2310.06474}, 2024.

\bibitem{yong2024lowresource}
Z.-X.~Yong, C.~Menghini, S.~H.~Bach,
Low-resource languages jailbreak GPT-4,
\textit{arXiv preprint arXiv:2310.02446}, 2024.

\bibitem{wei2022chain}
J.~Wei, X.~Wang, D.~Schuurmans, M.~Bosma, B.~Ichter, F.~Xia, E.~Chi,
Q.~V.~Le, D.~Zhou,
Chain-of-thought prompting elicits reasoning in large language models,
\textit{Advances in Neural Information Processing Systems},
\textbf{35}, 24824--24837, 2022.

\bibitem{kwon2023efficient}
W.~Kwon, Z.~Li, S.~Zhuang, Y.~Sheng, L.~Zheng, C.~H.~Yu,
J.~E.~Gonzalez, H.~Zhang, I.~Stoica,
Efficient memory management for large language model serving with
PagedAttention,
\textit{Proceedings of the 29th Symposium on Operating Systems Principles},
611--626, 2023.

\end{thebibliography}

% ============================================================
% APPENDIX  (after references, per journal style)
% ============================================================
\clearpage
\appendix

\section{Additional Method Details}
\label{sec:app-method}

\textbf{Refusal detection.}
We apply keyword-based refusal matching after removing any thinking trace
for reasoning models.
The pattern set includes common refusal markers (e.g., ``I cannot,''
``I'm sorry,'' ``as an AI'') and variants capturing hedged refusals.
This provides a transparent, reproducible baseline while avoiding
template-specific signals in the reasoning trace.

\textbf{Trace classification.}
We map each reasoning trace to Types~A--D using rule-based cues for explicit
safety awareness and the final response outcome.
Type~B is subdivided into explicit rationalizations (B-explicit) and
implicit compliance following an expressed concern (B-implicit), enabling
separation of ``self-persuasion'' from weak constraint propagation.

\textbf{Inference protocol.}
All models are evaluated with greedy decoding, max 512 new tokens, and
default chat templates.
This isolates the effect of training objective (instruction tuning vs.\
reasoning distillation) from sampling randomness or prompt engineering.

% ── Appendix: Extended analysis ──────────────────────────────
\section{Extended Analysis}
\label{sec:app-analysis}

\textbf{Cross-baseline robustness.}
Across all three architecture-matched pairs, reasoning distillation
consistently increases ASR relative to standard instruct training.
This holds despite large differences in baseline refusal rates
(95--99\% vs.\ 75--80\%), confirming that safety regression is not driven
by a single model family or parameter scale.

\textbf{Failure mode composition.}
Type~B and Type~C traces account for the majority of unsafe responses.
The split between these modes provides actionable mitigation targets:
Type~C suggests detection/activation gaps, whereas Type~B suggests
constraint propagation failures between reasoning and generation phases.

\textbf{Boundary-clarity hypothesis.}
The sharp contrast between deception (high failure) and
discrimination/sexual content (near-perfect refusal) supports a
boundary-clarity interpretation: categories with bright-line social norms
retain safety through distillation, while ambiguous helpful/harmful
boundaries become exploitable under deliberation.

\textbf{Generation length.}
Reasoning models produce substantially longer outputs (${\sim}$2{,}400 chars
vs.\ 250--870 chars), consistent with extended deliberation.
This increased generation budget may broaden the space of rationalizations
and weaken early refusal heuristics, offering a mechanism-level
interpretation for the observed ASR convergence.

% ── Appendix: Supplementary figures ──────────────────────────
\section{Supplementary Figures}
\label{sec:app-figs}

\begin{figure*}[!t]
\centering
\figwideh{figA1_paired_asr.png}
\caption{Paired ASR comparison across architecture-matched models.
Reasoning distillation consistently elevates ASR relative to the standard
instruct baseline.  Yellow badges show amplification factors.}
\label{fig:app-paired}
\end{figure*}

\begin{figure*}[!t]
\centering
\figwideh{figA2_radar.png}
\caption{Radar view of category-wise refusal rates (averaged across model
families).  The gap between the green (standard) and red (reasoning)
polygons visualizes category-selective degradation.}
\label{fig:app-radar}
\end{figure*}

\begin{figure*}[!t]
\centering
\figwideh{figA3_length.png}
\caption{Average response length by model.  Reasoning models consistently
generate 3--10$\times$ longer outputs than their standard counterparts,
reflecting extended deliberation.}
\label{fig:app-length}
\end{figure*}

\begin{figure*}[!t]
\centering
\figwideh{figA4_qualitative.png}
\caption{Representative thinking trace examples illustrating the three
safety-relevant reasoning patterns: Type~A (correct refusal after safety
awareness), Type~B (self-rationalization leading to compliance), and
Type~C (compliance without any safety awareness).}
\label{fig:app-qual}
\end{figure*}

\begin{figure*}[!t]
\centering
\figwideh{figA5_combined.png}
\caption{Combined overview of key quantitative trends:
(a)~ASR by model pair,
(b)~category-level safety gap,
(c)~trace type composition, and
(d)~response length comparison.}
\label{fig:app-combined}
\end{figure*}

\begin{figure*}[!t]
\centering
\figwideh{figA6_heatmap.png}
\caption{Full refusal-rate heatmap across all models (columns) and harm
categories (rows).  The white divider separates standard (left) and
reasoning (right) models.  Green = high refusal; red = low refusal.}
\label{fig:app-heatmap}
\end{figure*}

\end{document}
